%% Cross-volume comparison for the Calabi--Yau-to-chiral boundary.

\documentclass[10pt]{article}
\usepackage[localtheorems]{raeez-math-template}

\providecommand{\Vone}{\textsc{Vol~I}}
\providecommand{\Vtwo}{\textsc{Vol~II}}
\providecommand{\Vthree}{\textsc{Vol~III}}
\providecommand{\kch}{\kappa_{\mathrm{ch}}}
\providecommand{\kcat}{\kappa_{\mathrm{cat}}}
\providecommand{\kBKM}{\kappa_{\mathrm{BKM}}}
\providecommand{\kfib}{\kappa_{\mathrm{fiber}}}
\providecommand{\kHeis}{\kappa_{\mathrm{ch}}^{\mathrm{Heis}}}
\providecommand{\gDelta}{\mathfrak{g}_{\Delta_5}}
\providecommand{\HDelta}{\mathbf{H}_{\Delta_5}}
\providecommand{\CoHA}{\mathrm{CoHA}}
\providecommand{\Perf}{\mathrm{Perf}}
\providecommand{\SpCh}{\mathrm{Sp}^{\mathrm{ch}}}
\providecommand{\PhiFA}{\Phi^{\mathrm{FA}}}
\providecommand{\Eone}{E_1}
\providecommand{\Etwo}{E_2}
\providecommand{\cW}{\mathcal{W}}

\begin{document}

\begin{center}
\textbf{Cross-volume comparison for the Calabi--Yau-to-chiral boundary}\\[2pt]
\small The Vol~I, Vol~II, Vol~III, and Igusa manuscripts are used here as
fallible witnesses.  The displayed formulas are fixed by definitions,
direct Hodge or product computations, and named primary theorems.
\end{center}

\section*{Invariant separation on \(K3\times E\)}

Let \(Y=K3\times E\).  The compact Hodge column is
\[
        h^{0,\bullet}(Y)=(1,1,1,1),
\]
and therefore
\[
        \kch(Y)=\sum_q(-1)^q h^{0,q}(Y)=0.
\]
The categorical Euler characteristic is K\"unneth-multiplicative:
\[
        \kcat(Y)=\chi(\mathcal O_{K3})\chi(\mathcal O_E)=2\cdot0=0.
\]
The Heisenberg--Mukai specialization, the Borcherds weight, and the
K3 fibre rank are different constructions:
\[
 \{\kcat(Y),\,\kch(Y),\,\kHeis(Y),\,\kBKM(\gDelta),\,\kfib(Y)\}
 =
 \{0,\,0,\,3,\,5,\,24\}.
\]
The value \(3\) is the \(2+1\) Heisenberg--Mukai rank on the framed
\((K3,E)\) specialization.  The value \(5\) is the Borcherds weight of
Gritsenko's \(\Delta_5\).  The value \(24\) is the Mukai-lattice rank
of the K3 fibre.  These numbers cannot be combined by a sum law
unless a separate comparison morphism and its functoriality have first
been constructed.

\begin{proposition}[Separation of the four invariant formulas]
\label{prop:crossvol-no-additive-invariant-law}
For the primitive CHL denominators \(\Phi_N\), \(N\in\{1,2,3,4,6\}\),
the BKM invariant is
\[
        \kBKM(\Phi_N)=\frac{c_N(0)}{2}.
\]
It is the automorphic weight of the chosen Borcherds input.  It is
separate from \(\kch\), \(\kcat\), and \(\kfib\).
\end{proposition}

\begin{proof}
Borcherds' singular theta lift assigns to a weak Jacobi input with
constant term \(c_N(0)\) a product of weight \(c_N(0)/2\).  The
K\"unneth computation above gives \(\kcat(Y)=\kch(Y)=0\) on the compact
total space.  The K3 fibre contributes \(\kfib=24\), while
\(\kHeis=3\) comes from the framed Heisenberg specialization.  Since
these four constructions have different source functors, equality of
their numerical values has no invariant meaning without an explicit
comparison functor.  At \(N=1\), the Borcherds side gives
\(\kBKM(\Delta_5)=5\), while the compact total-space and elliptic-fibre
terms give \(0\) and \(0\).  Across
\(N\in\{1,2,3,4,6\}\), the Borcherds side gives
\((5,4,3,2,1)\).  No Hodge-supertrace or fibre-Euler construction
has these values.  The obstruction is structural, and the formula
\(\kBKM(\Phi_N)=c_N(0)/2\) is the surviving invariant.
\end{proof}

\section*{Primitive CHL denominators and Igusa normalization}

The primitive smooth \(K3\times E\) CHL denominator row is indexed by
\[
        N=(1,2,3,4,6).
\]
With the Borcherds input normalized at the relevant cusp, the constants
and weights are
\[
\begin{array}{c|ccccc}
N & 1 & 2 & 3 & 4 & 6\\
\hline
c_N(0) & 10 & 8 & 6 & 4 & 2\\
\kBKM(\Phi_N)=c_N(0)/2 & 5 & 4 & 3 & 2 & 1 .
\end{array}
\]
At level one,
\[
        \Delta_5=\operatorname{Borch}(\phi_{0,1}),\qquad
        \phi_{0,1}(\tau,z)=r^{-1}+10+r+O(q),
\]
so \(c_1(0)=10\) and \(\mathrm{wt}(\Delta_5)=5\).  The monic
denominator product is
\[
        D_5=64^{-1}\Delta_5,\qquad
        \operatorname{den}(\gDelta)=D_5(2Z)=64^{-1}\Delta_5(2Z).
\]
The Igusa square and the reduced \(K3\times E\) scalar use different
normalizations:
\[
        \Phi_{10}^{\mathrm{un}}=\Delta_5^2,\qquad
        \Phi_{10}^{\mathrm{OP}}=D_5^2=4096^{-1}\Delta_5^2,
\]
\[
        Z^{K3\times E}_{\mathrm{OP/DT}}
        =-(\Phi_{10}^{\mathrm{OP}})^{-1}
        =-4096\,\Delta_5^{-2}.
\]
The K3 elliptic genus is \(2\phi_{0,1}\).  Its Borcherds lift has
weight \(10\) and lands on \(\Phi_{10}^{\mathrm{un}}\).  The primitive
BKM denominator uses \(\phi_{0,1}\), weight \(5\), and the square-root
normalization above.
The factor \(64\) is the theta-leading coefficient of \(\Delta_5\).
The factor \(4096\) is its square in the Oberdieck--Pandharipande
scalar convention.  Neither factor changes the BKM weight \(5\).

The Clery--Gritsenko diagonal-divisor catalogue has eight rows, with
weights
\[
        (5,2,3,1,2,\tfrac12,\tfrac32,1)
\]
and product constants
\[
        (10,4,6,2,4,1,3,2).
\]
This table is a diagonal-divisor catalogue, while
\((10,8,6,4,2)\) is the primitive CHL denominator row.  A comparison
between them requires a row map, a group character, a multiplier cover,
and a Borcherds root datum.

The orders \(N=5,7,8\) belong to Nikulin's K3 symplectic-automorphism
boundary.  A denominator theorem at these orders must specify
\[
        X_N,\quad L_N,\quad \rho_N,\quad
        \Delta_N^+,\quad \operatorname{mult}_N,\quad
        F^{(g_N,h_M)}_{L_N}.
\]
Once those data exist, Borcherds' theorem again gives
\(\kBKM(F^{(g_N,h_M)}_{L_N})=c_N(0)/2\).  The compact
\(K3\times E\) CHL quotient geometry supplies the primitive row only at
\(N\in\{1,2,3,4,6\}\).

\section*{Operadic scope of \(\Phi_d\)}

The dimension-\(d\) construction factors as
\[
        \Phi_d=\SpCh_{\Sigma_{d-1},C}\circ\PhiFA_d.
\]
Stage \(1\) produces the holomorphic factorization algebra attached to
the \(d\)-Calabi--Yau category.  Stage \(2\) integrates over
\(\Sigma_{d-1}\) and restricts to the reference curve \(C\).  For
\(d\geq3\), this operation leaves a native \(\Eone\)-chiral object:
\[
        \SpCh_{\Sigma_{d-1},C}
        \bigl(\PhiFA_d(\Perf(X))\bigr)\in
        \mathrm{ChirAlg}_{\Eone}(C).
\]
Braiding at \(\Etwo\)-level is recovered, on constructed loci, from
the Drinfeld centre, double, or representation category:
\[
        \mathcal Z\bigl(\mathrm{Rep}^{\Eone}(A)\bigr).
\]
Thus at \(d=3\) the object \(A\) is \(\Eone\)-native; the
\(\Etwo\)-structure belongs to the centre/double/module layer.

\section*{Positive halves and vertex targets}

For the toric point \(\mathbb C^3\), the Hall-side computation is
\[
        \CoHA(\mathbb C^3)\simeq Y^+(\widehat{\mathfrak{gl}}_1)
\]
by Schiffmann--Vasserot.  This is the positive half of the affine
Yangian and an \(\Eone\)-chiral Hall output.  The full
\(\cW_{1+\infty}\) object at \(c=1\) appears after the passage
\[
        Y^+(\widehat{\mathfrak{gl}}_1)
        \hookrightarrow
        Y(\widehat{\mathfrak{gl}}_1)
        \xrightarrow{\mathrm{ev}_\lambda}
        \operatorname{End}\bigl(\cW_{1+\infty}[\lambda]\text{-vacuum}\bigr),
\]
or equivalently through the Drinfeld centre/double and evaluation
module.  The bare Hall algebra has an ordered multiplication from
short-exact-sequence correspondences.  The \(R\)-matrix and the
non-symmetric braiding require the doubled or centred layer.

\section*{Four branches}

The comparison splits into four branches.
\begin{center}
\renewcommand{\arraystretch}{1.18}
\begin{tabular}{p{0.22\linewidth}p{0.31\linewidth}p{0.35\linewidth}}
\toprule
\textbf{Branch} & \textbf{Input} & \textbf{Output}\\
\midrule
BKM / Hall--Drinfeld &
K3-fibred \(K3\times E\), Borcherds input, Hall positive half &
\(\gDelta\), \(\HDelta\), and \(D(Y^+_{\mathrm{Hall}}(K3\times E))\) on
constructed principal loci\\
Yangian / Mukai &
K3 Mukai lattice, self-mirror and stable-envelope data &
K3 Yangian-type or toroidal objects governed by Mukai geometry\\
Class-\(\mathcal S\) / VOA &
Schur sector, \(A_1\) class-\(\mathcal S\) curve \(\Sigma_{0,24}\),
corner-VOA data &
VOA and chiral-bialgebra avatars carrying their own representation
categories\\
DT / GW / MNOP &
Reduced DT, Gromov--Witten, BCOV, and topological vertex data &
Partition functions and BPS counts; BKM recognition only after a
Borcherds root datum is supplied\\
\bottomrule
\end{tabular}
\end{center}
The branches can meet through explicit functors, centres, doubles,
stable envelopes, or partition-function identities.  A shared number or
symbol alone is insufficient.

\section*{Fifteen comparison loci}

\begin{enumerate}
\item \emph{Non-abelian \(6d\) hCS on the quintic.}
The abelian and reductive hCS formalisms supply the local BV
factorization-algebra input.  A compact quintic theorem also requires
the non-formal \(\mathbb S^3\)-framing datum, compact BV quantization,
and the oriented hCS-to-Hall comparison.  The BCOV/MNOP branch is the
natural compact-CY\(_3\) replacement for a Borcherds denominator.

\item \emph{Primitive CHL denominators at \(N=2,3,4,6\).}
The automorphic weights are fixed by
\[
        c_N(0)=(8,6,4,2),\qquad
        \kBKM(\Phi_N)=(4,3,2,1).
\]
The enumerative realization uses the reduced DT geometry of the
corresponding \(K3\times E\) quotient and an equivariant Hall--Borcherds
comparison.  The automorphic theorem supplies the weight; the
geometry supplies the realization.

\item \emph{Boundary orders \(N=5,7,8\).}
These orders carry K3 twining arithmetic and multiplier-cover
automorphic forms.  Compact \(K3\times E\) CHL
geometry gives no primitive row at these orders.  A theorem needs a
host \(X_N\), a root lattice, a Weyl chamber, and a multiplicity
function.

\item \emph{MNOP for generic compact CY\(_3\).}
For \(\mathbb C^3\), the bar complex of \(Y^+(\widehat{\mathfrak{gl}}_1)\)
recovers the MacMahon function and the DT invariants
\(\Omega_{\mathrm{DT}}(n)=n\).  For a generic compact CY\(_3\), MNOP
belongs to the DT/GW/BCOV branch.  Borcherds recognition requires an
additional lattice and denominator datum.

\item \emph{Chiral Booth--Lazarev.}
The curved Quillen and coderived model structures give the abstract
homotopical machine for curved coalgebras and modules.  A chiral
instantiation must exhibit the corresponding factorization-algebra
objects, monoidal model structure, and compatibility with the
\(\Eone\)-chiral envelope.

\item \emph{Orthosymplectic super-Yangian.}
The Mukai lattice has signature \((4,20)\), and the structure-preserving
classical algebra is orthosymplectic.  Thus the conjectural object is
\[
        Y_{\mathfrak{osp}(4|20)}
\]
with reflection-equation and Berezinian data.  The classical
\(\mathfrak{gl}(m|n)\) super-Yangian literature supplies comparison
technology; the K3 object requires the orthosymplectic form.

\item \emph{CY\(_4\) prediction.}
The compact Hodge supertrace formula remains
\[
        \kch(A_X)=\sum_q(-1)^q h^{0,q}(X).
\]
The \(d=4\) output is \(E_1\)-stabilized after Stage \(2\), with any
braided structure carried by a centre or module category.  A complete
\(\Phi_4\) theorem requires the relevant framing and functoriality
data.

\item \emph{\(G(X)\) for generic \(X\).}
For a constructed positive half \(Y^+(X)\) with nondegenerate Hopf
pairing and completed tensor product, one may define
\[
        G(X)=D(Y^+(X)).
\]
For \(K3\times E\), the six routes to \(G(K3\times E)\) are six
constructions of the same principal object: Hall/SV CoHA,
Frenkel--Kac/Borcherds imaginary roots, Siegel-paramodular
\(R\)-matrix, bi-based Ran factorization, Costello one-loop hCS, and
class-\(\mathcal S\) chiral algebra.  They are separate constructions,
all requiring their own hypotheses.

\item \emph{CY-C on generic K3.}
The expected comparison
\[
        C(\mathfrak g,q)=D\bigl(Y^+(\mathfrak g_{K3})\bigr)
\]
depends on a K3 positive half, a Hopf pairing, stable-envelope
compatibility, and centre/double passage.  The class-\(\mathcal S\)
Schur-sector realization gives compatible VOA evidence, while the
Hall--Drinfeld construction supplies the algebraic target.

\item \emph{Class \(\mathcal M\) at high genus.}
The cohomological class-\(\mathcal M\) \(E_3\)-bar count is \(6^g\).
The raw direct-sum chain model already fails at \(g=2\).  The correct
ambient is pro-completed, \(J\)-adic, or filtered-completed.  At
\(g\ge4\), the Schottky and Humbert boundary components add geometric
constraints to the same completed setting.

\item \emph{Half-integral diagonal-divisor rows.}
The Clery--Gritsenko rows of weights \(1/2\) and \(3/2\) live on
multiplier covers.  Their product constants are catalogue-row
constants, distinct from the primitive CHL tuple.  A physical or
enumerative interpretation must name the cover, the Jacobi input, and
the denominator algebra.

\item \emph{Fake Monster sibling.}
Borcherds' Fake Monster Lie algebra is attached to
\(II_{25,1}\) and the Leech lattice.  The compact \(d=3\) placement is
blocked by rank and parity.  The natural CY-to-chiral routing is a
conditional \(d=5\) construction on \(K3\times K3\times E\), using
\[
        E_5\simeq E_2\otimes E_2\otimes E_1.
\]

\item \emph{Humbert surface sibling.}
The \(\mathcal B\)-family carries
\[
        K^{\kch}=8=\operatorname{ord}(H_1),\qquad
        \hbar^2 K^{\kch}=-1.
\]
The Humbert strata \(H_1\cup H_4\) support the gerbe obstruction for
the \(\Delta_5\) comparison, while \(H_1\cup H_3\cup H_4\) supports the
chain-level Malcev-ladder inversion in the completed class-\(\mathcal M\)
setting.

\item \emph{Central value of the Igusa square.}
The spin \(L\)-function factorization is
\[
        L(s,\Phi_{10}^{\mathrm{un}})
        =
        L(s,\Delta\otimes E_6)\zeta(s-9)\zeta(s-8).
\]
The value \(L(5/2,\Phi_{10}^{-1})\) requires an analytic normalization
of the inverse Igusa square and a regularized central-value statement.
The factorization alone supplies the \(L\)-function line; the
central-value evaluation requires the inverse-square normalization and
regularization.

\item \emph{\(\gDelta\) to \(\HDelta\).}
The Lie algebra \(\gDelta\) is the Borcherds--Kac--Moody algebra with
denominator \(D_5(2Z)\).  The chiral bialgebra \(\HDelta\) is the
Hall--Drinfeld double with Gritsenko--Nikulin denominator cocycle
\(m_0=\omega_{\Delta_5}\), after the positive half, pairing,
completion, and Borcherds comparison have been supplied.  The Mukai
self-mirror K3 Yangian branch remains a separate comparison branch.
\end{enumerate}

\section*{Source spine}

The local source spine consists of the Vol~III \(\kappa_\bullet\)
stratification chapter, the Vol~III \(\Phi_d\) construction chapter,
the two \(\Delta_5\) automorphic notes, the Vol~I landscape census,
the Vol~II doctrine file, and the Igusa manuscript.
The primary source spine is Borcherds' singular theta lift and
generalized Kac--Moody theory, Gritsenko--Nikulin's \(\Delta_5\)
denominator construction, Clery--Gritsenko's diagonal-divisor
classification, Schiffmann--Vasserot's CoHA identification for
\(\mathbb C^3\), Kontsevich--Soibelman cohomological Hall algebras,
Costello--Gwilliam factorization algebras, and
Oberdieck--Pandharipande/Oberdieck--Pixton reduced \(K3\times E\)
Donaldson--Thomas theory.

\end{document}
